\documentclass{report} % ou book
\usepackage[T1]{fontenc}      
\usepackage[utf8]{inputenc}    
\usepackage{lmodern}           
\usepackage{amsmath, amssymb}  
\usepackage{mathtools}         
\usepackage{enumitem}          
\usepackage{geometry}
\usepackage{amsthm}  % Pour définir theorems, definitions, etc.

% Définition des environnements
\theoremstyle{definition}  % style romain, pas en italique
\newtheorem{definition}{Définition}[chapter] % numérotation par chapitre
\newtheorem{example}{Exemple}[chapter]
\newtheorem{corollary}{Corollaire}[chapter] % <- pour les corollaires
\theoremstyle{plain} % pour propositions et théorèmes
\newtheorem{theorem}{Théorème}[chapter]
\newtheorem{proposition}{Proposition}[chapter]
\newtheorem{axiom}{Axiome}
\theoremstyle{remark} % pour les remarques
\newtheorem*{remark}{Remarque} % pas numérotée

\geometry{margin=2.5cm}

\title{Analyse : Mesures et intégration (Partie I)}
\author{David Thébault}
\date{\today}

\begin{document}

\maketitle

\chapter{\textbf{Intégrales}}

\section{\textbf{Rappels}}

\subsection{Notions de base}

Soit \(f : [a,b] \subset \mathbb{R} \to \mathbb{R}\) une fonction continue.  
L'intégrale de \(f(x)\,dx\) sur \([a,b]\) représente, si \(a < b\), l'aire algébrique du domaine délimité par la courbe \(C_f\) de \(f\) et les droites \(y=0\), \(x=a\) et \(x=b\) :  

- la partie située au-dessus de l'axe des abscisses contribue positivement,  
- la partie située en dessous de l'axe des abscisses contribue négativement.  

\[
\int_a^b f(x)\,dx
\]

\noindent
La lettre \(x\) est dite "muette" : on peut écrire  
\[
\int_a^b f(x)\,dx = \int_a^b f(y)\,dy
\]  
et on peut noter également simplement  
\[
\int_a^b f \quad \text{ou} \quad \int f \text{ entre } a \text{ et } b.
\]

\medskip
\noindent
Si \(a > b\), on a la convention :  
\[
\int_a^b f = - \int_b^a f
\]

\medskip

\subsection{Intégrabilité au sens de Riemann}

Soit \(f : [a,b] \to \mathbb{R}\) une fonction bornée.
On approche l’aire algébrique sous la courbe représentative de \(f\)
par des \emph{sommes de Riemann}.
Pour une subdivision
\[
a = x_0 < x_1 < x_2 < \dots < x_n = b,
\]
on associe à chaque sous-intervalle \([x_i,x_{i+1}]\) un rectangle
de base \(x_{i+1}-x_i\) et de hauteur \(f(\xi_i)\), où le point
\(\xi_i\) est choisi dans \([x_i,x_{i+1}]\).

\medskip
\noindent
Le point \(\xi_i\) peut être choisi arbitrairement dans l’intervalle
\([x_i,x_{i+1}]\).
Un choix fréquent est le point milieu
\[
\xi_i = \frac{x_i + x_{i+1}}{2}.
\]
Lorsque le pas de la subdivision tend vers 0, la limite des sommes de Riemann,
lorsqu’elle existe, ne dépend pas du choix des points \(\xi_i\) et définit
l’intégrale de \(f\) sur \([a,b]\).


\medskip
\noindent
La fonction \(f\) est dite \emph{intégrable au sens de Riemann} sur \([a,b]\)
si ces sommes convergent vers une même limite lorsque le pas de la subdivision
tend vers 0. Cette limite est notée
\[
\int_a^b f(x)\,dx.
\]

\medskip


\begin{definition}[Intégrabilité au sens de Riemann]
Soit \( f : [a, b] \to \mathbb{R} \) une fonction bornée. On dit que \( f \) est \textbf{intégrable au sens de Riemann} sur \([a, b]\) si pour toute subdivision \(\sigma\) de \([a, b]\) définie par :
\[
\sigma = \{x_0 = a < x_1 < x_2 < \dots < x_n = b\},
\]
et pour tout choix de points intermédiaires \(\xi_i \in [x_{i-1}, x_i]\), les \textbf{sommes de Riemann} :
\[
S(\sigma, f) = \sum_{i=1}^n f(\xi_i)(x_i - x_{i-1}),
\]
convergent vers une même limite \(I\) lorsque le \textbf{pas de la subdivision} \(\|\sigma\| = \max_{1 \leq i \leq n} (x_i - x_{i-1})\) tend vers 0.

\medskip
\noindent
Cette limite commune \(I\) est appelée \textbf{l'intégrale de Riemann} de \(f\) sur \([a, b]\), et elle est notée :
\[
\int_a^b f(x) \, dx.
\]
\end{definition}

\medskip

\textbf{Remarques :}
\begin{itemize}
    \item Le pas de la subdivision \(\|\sigma\|\) est la longueur du plus grand sous-intervalle \([x_{i-1}, x_i]\).
    \item Si \(f\) est continue sur \([a, b]\), alors elle est intégrable au sens de Riemann.
    \item Les sommes de Riemann \(S(\sigma, f)\) approchent l'aire sous la courbe de \(f\) par des rectangles de hauteur \(f(\xi_i)\) et de largeur \(x_i - x_{i-1}\).
\end{itemize}

\medskip

\begin{theorem}[\textbf{Critère d’intégrabilité de Riemann}]
Une fonction \(f : [a,b] \to \mathbb{R}\) est intégrable au sens de Riemann
si et seulement si :
\begin{itemize}
    \item \(f\) est \textbf{bornée} sur \([a,b]\),
    \item l’ensemble de ses \textbf{points de discontinuité est de mesure nulle}
    (pour la mesure de Lebesgue), c’est-à-dire qu’il est « négligeable »,
    par exemple un ensemble fini ou dénombrable de points.
\end{itemize}
\end{theorem}

\medskip
\noindent
Exemples de fonctions intégrables au sens de Riemann :  

\medskip
La fonction en escalier définie par:
\[
f(x) =
\begin{cases}
1 & 0 \le x < 1,\\
2 & 1 \le x \le 2
\end{cases}
\]
est discontinue en un nombre fini de points mais intégrable sur \([0,2]\).

\medskip
La fonction composée \(f\) définie par:
\[
f(x) = \begin{cases} \sin(1/x) & x \in (0,1], \\ 0 & x = 0 \end{cases}
\]
est discontinue en \(0\) mais bornée et intégrable sur \([0,1]\).

\medskip
La fonction indicatrice de l’ensemble de Cantor 
\(\chi_C : [0,1] \to [0,1]\) définie par
\[
\chi_C(x) =
\begin{cases}
1 & x \in C, \\
0 & x \notin C,
\end{cases}
\]
où \(C\) est l’ensemble de Cantor construit en retirant à chaque étape le tiers central des intervalles restants.  

\noindent
Cette fonction est bornée et discontinue uniquement sur l’ensemble \(C\),  
qui est de mesure de Lebesgue nulle. Ainsi, \(\chi_C\) est Riemann intégrable.
\medskip
\begin{remark}[Mesure nulle de l'ensemble de Cantor]
Pour construire l'ensemble de Cantor \(C\), on commence avec l'intervalle \([0,1]\) 
et on retire le tiers central \((1/3, 2/3)\), laissant les intervalles \([0,1/3]\) et \([2/3,1]\).
Ensuite, on retire le tiers central de chaque intervalle restant, 
c'est-à-dire \((\frac{1}{9}, \frac{2}{9})\) et \((\frac{7}{9}, \frac{8}{9})\),
laissant quatre intervalles : \([0,1/9]\), \([2/9,1/3]\), \([2/3,7/9]\), et \([8/9,1]\).
On répète indéfiniment le processus sur chaque intervalle restant : à chaque étape \(n\), il reste \(2^n\) intervalles de longueur \((1/3)^n\).  

On peut utiliser ces intervalles comme recouvrement de \(C\) :
\[
C \subset \bigcup_{k=1}^{2^n} I_k, \quad \ell(I_k) = \left(\frac{1}{3}\right)^n.
\]
La somme des longueurs est
\[
\sum_{k=1}^{2^n} \ell(I_k) = 2^n \cdot \left(\frac{1}{3}\right)^n = \left(\frac{2}{3}\right)^n.
\]
Pour tout \(\varepsilon > 0\), il existe \(n\) assez grand pour que
\((2/3)^n < \varepsilon\). Ainsi, \(C\) peut être recouvert par une union finie
d’intervalles dont la somme des longueurs est inférieure à \(\varepsilon\),
ce qui montre que \(C\) est de mesure de Lebesgue nulle, bien qu’il soit non dénombrable.
\end{remark}

\medskip
\noindent
Exemple de fonction non intégrable au sens de Riemann :  

\medskip
Exemple : la fonction de Dirichlet \(f_D : [0,1] \to \mathbb{R}\) définie par
\[
f_D(x) =
\begin{cases}
1 & x \in \mathbb{Q},\\
0 & x \notin \mathbb{Q}.
\end{cases}
\]
Elle est discontinue en tout point de \([0,1]\) et donc non intégrable au sens de Riemann.




\subsection{Mesure de Lebesgue (introduction)}

La mesure de Lebesgue sur \(\mathbb{R}\) généralise la notion de longueur
des intervalles.
Ainsi, pour tout intervalle \([a,b]\), sa mesure de Lebesgue est
\[
\lambda([a,b]) = b-a.
\]

De manière intuitive, la mesure de Lebesgue d’un ensemble de réels
représente sa « taille » ou sa « longueur » sur l’axe réel.


\medskip
\begin{definition}[Mesure de Lebesgue pour un ensemble]
Soit $E \subset \mathbb{R}$. On dit que $E$ est \emph{mesurable au sens de Lebesgue} et sa \emph{mesure de Lebesgue} $m(E)$ est définie par :
\[
m(E) = \inf \Bigg\{ \sum_{n=1}^{\infty} \ell(I_n) \;\Big|\; E \subset \bigcup_{n=1}^{\infty} I_n \Bigg\},
\]
où $(I_n)_{n \ge 1}$ est une suite dénombrable d'intervalles ouverts et $\ell(I_n)$ désigne la longueur de l'intervalle $I_n$.

Autrement dit, $m(E)$ est la borne inférieure des sommes des longueurs des recouvrements dénombrables d'intervalles qui contiennent $E$.
\end{definition}


\medskip
Exemple : tout intervalle $[a,b]$ est mesurable et $m([a,b]) = b-a$.  
Tout ensemble fini ou dénombrable de points est mesurable et de mesure $0$.


\medskip
\begin{definition}[Mesure de Lebesgue nulle]
Un ensemble \(E \subset \mathbb{R}\) est dit de \emph{mesure de Lebesgue nulle} si, pour tout \(\varepsilon > 0\), il existe une suite d'intervalles ouverts 
\((I_n)_{n \ge 1}\) tels que :
\[
E \subset \bigcup_{n=1}^{\infty} I_n
\quad \text{et} \quad
\sum_{n=1}^{\infty} \ell(I_n) \le \varepsilon,
\]
où \(\ell(I_n)\) désigne la longueur de l'intervalle \(I_n\).
\end{definition}

\medskip
\noindent
Un ensemble est dit de \emph{mesure de Lebesgue nulle} s’il peut être recouvert,
pour tout \(\varepsilon > 0\), par une union dénombrables d’intervalles dont la somme des
longueurs est inférieure à \(\varepsilon\).
Autrement dit, un tel ensemble est négligeable du point de vue de la longueur.

\medskip
\begin{remark}
Tout ensemble fini ou dénombrable de points est de mesure de Lebesgue nulle.
Ainsi, toute fonction bornée et continue sauf en un nombre fini de points
est intégrable au sens de Riemann.
\end{remark}
\medskip
\noindent
Exemple 1 : ensemble fini

Soit \(E = \{2,3\} \subset \mathbb{R}\). Montrons que \(E\) est de mesure de Lebesgue nulle.

Soit \(\varepsilon > 0\). On considère les intervalles suivants recouvrant chaque point de \(E\) :
\[
I_1 = \left(2 - \frac{\varepsilon}{4}, 2 + \frac{\varepsilon}{4}\right), \quad
I_2 = \left(3 - \frac{\varepsilon}{4}, 3 + \frac{\varepsilon}{4}\right).
\]
Alors \(E \subset I_1 \cup I_2\) et la longueur totale est
\[
\ell(I_1) + \ell(I_2) = \frac{\varepsilon}{2} + \frac{\varepsilon}{2} = \varepsilon.
\]
Ainsi \(E\) est de mesure de Lebesgue nulle.

\medskip
\noindent
Exemple 2 : ensemble dénombrable

Soit \(E = \{x_1, x_2, x_3, \dots\} \subset \mathbb{R}\) un ensemble dénombrable.
Pour tout \(\varepsilon > 0\), recouvrons chaque point \(x_n\) par l'intervalle
\[
I_n = \left(x_n - \frac{\varepsilon}{2^{n+1}}, \, x_n + \frac{\varepsilon}{2^{n+1}}\right).
\]
Alors :
\[
E \subset \bigcup_{n=1}^{\infty} I_n, \quad 
\sum_{n=1}^{\infty} \ell(I_n) = \sum_{n=1}^{\infty} \frac{\varepsilon}{2^n} = \varepsilon.
\]
Ainsi tout ensemble dénombrable de points est de mesure de Lebesgue nulle.



\section{\textbf{Propriétés des intégrales}}

Soit \(f, g : [a,b] \to \mathbb{R}\) deux fonctions continues, et \(\lambda \in \mathbb{R}\).  

\subsection{Linéarité}

\[
\int_a^b (\lambda f(x) + g(x)) \, dx = \lambda \int_a^b f(x) \, dx + \int_a^b g(x) \, dx
\]

\noindent
\textbf{Preuve :}  
On utilise la définition de l'intégrale comme limite de sommes de Riemann. Pour un découpage \(a = x_0 < x_1 < \dots < x_n = b\) et \(\xi_i \in [x_{i-1},x_i]\) :

\[
\sum_{i=1}^{n} (\lambda f(\xi_i) + g(\xi_i)) \Delta x_i = \lambda \sum_{i=1}^{n} f(\xi_i) \Delta x_i + \sum_{i=1}^{n} g(\xi_i) \Delta x_i
\]

En passant à la limite lorsque la taille des intervalles tend vers 0, on obtient la linéarité.

\subsection{Relation de Chasles}

\[
\int_a^b f(x) \, dx = \int_a^c f(x) \, dx + \int_c^b f(x) \, dx, \quad \forall c \in [a,b]
\]

\noindent
\textbf{Preuve :}  
On découpe l’intégrale en deux sous-intervalles \([a,c]\) et \([c,b]\).  
La somme de Riemann de \([a,b]\) peut se réécrire comme la somme des sommes de Riemann sur \([a,c]\) et \([c,b]\).  
En passant à la limite, on obtient la relation de Chasles.

\subsection{Positivité}

Si \(f(x) \ge 0\) pour tout \(x \in [a,b]\), alors

\[
\int_a^b f(x) \, dx \ge 0
\]

\noindent
\textbf{Preuve :}  
Chaque terme de la somme de Riemann \(\sum f(\xi_i) \Delta x_i \ge 0\).  
La limite (l’intégrale) est donc \(\ge 0\).

\subsection{Comparaison}

Si \(f(x) \le g(x)\) pour tout \(x \in [a,b]\), alors

\[
\int_a^b f(x) \, dx \le \int_a^b g(x) \, dx
\]

\noindent
\textbf{Preuve :}  
Considérons \(h(x) = g(x) - f(x) \ge 0\) pour tout \(x \in [a,b]\).  
Alors, par positivité :

\[
\int_a^b h(x) \, dx \ge 0 \implies \int_a^b (g(x) - f(x)) \, dx \ge 0
\]

\noindent
ce qui donne

\[
\int_a^b g(x) \, dx - \int_a^b f(x) \, dx \ge 0 \implies \int_a^b f(x) \, dx \le \int_a^b g(x) \, dx
\]

\subsection{Inégalité triangulaire}

Pour toute fonction \(f\) intégrable sur \([a,b]\) :

\[
\left| \int_a^b f(x)\,dx \right| \le \int_a^b |f(x)|\,dx
\]

\noindent
\textbf{Preuve :}  
Posons \(I = \int_a^b f(x)\,dx\). Alors, pour tout \(x \in [a,b]\), \(-|f(x)| \le f(x) \le |f(x)|\).  
En intégrant ces inégalités et en utilisant la propriété de comparaison, on obtient :

\[
-\int_a^b |f(x)|\,dx \le \int_a^b f(x)\,dx \le \int_a^b |f(x)|\,dx
\]

\noindent
ce qui implique

\[
\left| \int_a^b f(x)\,dx \right| \le \int_a^b |f(x)|\,dx
\]

\medskip

\begin{theorem}[\textbf{Théorème de la moyenne} pour les intégrales]
Si \( f \) est \textbf{continue} sur \([a, b]\), alors il existe \(\xi \in [a, b]\) tel que
\[
\int_a^b f(x) \, dx = f(\xi)(b - a).
\]
\end{theorem}

\medskip

\noindent
Autrement dit, l’intégrale de \( f \) sur \([a, b]\) est égale à la valeur de \( f \) en un certain point \(\xi\) multipliée par la longueur de l’intervalle \((b - a)\).

\medskip

\begin{proof}[Preuve]
Puisque \( f \) est continue sur \([a, b]\), elle atteint ses bornes inférieure et supérieure sur cet intervalle. Notons :
\[
m = \min_{x \in [a, b]} f(x) \quad \text{et} \quad M = \max_{x \in [a, b]} f(x).
\]
Pour tout \( x \in [a, b] \), on a \( m \leq f(x) \leq M \). En intégrant ces inégalités, on obtient :
\[
m(b - a) \leq \int_a^b f(x) \, dx \leq M(b - a).
\]
Divisons par \((b - a)\) (qui est positif) :
\[
m \leq \frac{1}{b - a} \int_a^b f(x) \, dx \leq M.
\]
Par le \textbf{théorème des valeurs intermédiaires}, il existe \(\xi \in [a, b]\) tel que :
\[
f(\xi) = \frac{1}{b - a} \int_a^b f(x) \, dx.
\]
En multipliant par \((b - a)\), on obtient :
\[
\int_a^b f(x) \, dx = f(\xi)(b - a).
\]
\end{proof}

\medskip

\textbf{Remarques :}
\begin{itemize}
    \item Ce théorème est souvent utilisé pour définir la \textbf{valeur moyenne} d’une fonction sur un intervalle : \(\frac{1}{b - a} \int_a^b f(x) \, dx\).
    \item Il repose sur la continuité de \( f \) sur un intervalle fermé \([a, b]\), ce qui garantit l’existence des valeurs extrêmes \( m \) et \( M \).
    \item Ce résultat est fondamental en analyse et en probabilités, où il permet de relier les notions d’intégrale et de moyenne.
    \item Il peut être généralisé à des fonctions à valeurs vectorielles ou à des intégrales multiples.
    \item Il est un outil clé pour établir des bornes sur les intégrales et pour prouver d’autres inégalités.
\end{itemize}


\medskip


\section{\textbf{Lien primitive – intégrale}}

\begin{definition}[Primitive]
Soit \(f\) une fonction définie sur un intervalle \(I \subset \mathbb{R}\).  
On appelle \emph{primitive} de \(f\) sur \(I\) toute fonction \(F\) dérivable sur \(I\) telle que
\[
F'(x) = f(x), \quad \forall x \in I.
\]
\end{definition}
\medskip
\begin{theorem}[Théorème fondamental du calcul intégral - Isaac Newton]
Soit \(f : I \to \mathbb{R}\) continue sur un intervalle \(I\) et \(a \in I\) fixé.  
Alors la fonction
\[
F : I \to \mathbb{R}, \quad F(x) = \int_a^x f(t)\, dt
\]
est de classe \(\mathcal{C}^1\) sur \(I\) et sa dérivée est
\[
F'(x) = f(x), \quad \forall x \in I.
\]
\end{theorem}
\medskip
\begin{proof}[Preuve]
Pour \(x \in I\) et \(h\) petit tel que \(x+h \in I\), on a :
\[
\frac{F(x+h) - F(x)}{h} = \frac{1}{h} \int_x^{x+h} f(t)\, dt.
\]
Par le théorème des accroissements finis pour les intégrales, il existe \(\xi \in [x, x+h]\) tel que
\[
\int_x^{x+h} f(t)\, dt = f(\xi) \cdot h.
\]
Donc :
\[
\frac{F(x+h) - F(x)}{h} = f(\xi).
\]
En faisant tendre \(h \to 0\), \(\xi \to x\) et par continuité de \(f\), on obtient :
\[
\lim_{h \to 0} \frac{F(x+h) - F(x)}{h} = f(x),
\]
ce qui montre que \(F\) est dérivable et que \(F' = f\).
\end{proof}
\medskip

\begin{corollary}
Soit \(f : [a,b] \to \mathbb{R}\) continue et \(F\) une primitive de \(f\) sur \([a,b]\).  
Alors :
\[
\int_a^b f(x)\, dx = F(b) - F(a),
\]
ce que l’on note également \([F(x)]_a^b\).
\end{corollary}
\medskip
\begin{proof}[Preuve]
Puisque \(F\) est une primitive de \(f\), on a \(F'(x) = f(x)\) pour tout \(x \in [a,b]\).  
D'après le théorème fondamental :
\[
\int_a^b f(x)\, dx = \int_a^b F'(x)\, dx = F(b) - F(a).
\]
\end{proof}

\begin{remark}
Ce corollaire permet de calculer facilement des intégrales définies en trouvant simplement une primitive de la fonction intégrande.
\end{remark}

\noindent
Exemple d'application du changement de variable
Calculons :
\[
\int_0^1 \sqrt{1-x^2}\, dx
\]

\noindent
On effectue le changement de variable \(x = \cos(t)\), donc \(dx = -\sin(t)\, dt\).  
Lorsque \(x = 0\), \(t = \pi/2\) ; lorsque \(x = 1\), \(t = 0\).  

Ainsi :
\[
\int_0^1 \sqrt{1-x^2}\, dx
= \int_{\pi/2}^{0} \sqrt{1-\cos^2(t)}\, (-\sin(t))\, dt
= \int_0^{\pi/2} \sin^2(t)\, dt.
\]

En utilisant l'identité \(\sin^2(t) = \frac{1 - \cos(2t)}{2}\), on obtient :
\[
\int_0^{\pi/2} \sin^2(t)\, dt
= \int_0^{\pi/2} \frac{1 - \cos(2t)}{2}\, dt
= \frac{1}{2} \left[ t - \frac{\sin(2t)}{2} \right]_0^{\pi/2}
= \frac{\pi}{4}.
\]
\medskip

\begin{theorem}
[Intégration par parties]

\end{theorem}
Soient \(u, v : [a,b] \to \mathbb{R}\) deux fonctions de classe \(\mathcal{C}^1\).
Alors :
\[
\int_a^b u(x)v'(x)\, dx = [u(x)v(x)]_a^b - \int_a^b u'(x)v(x)\, dx.
\]
\medskip

\begin{proof}[Preuve]
On utilise la formule de la dérivée du produit vue précédemment :
\[(uv)'(x) = u'(x)v(x) + u(x)v'(x).\]
En intégrant cette égalité sur \([a,b]\), on obtient :
\[
\int_a^b (uv)'(x)\, dx = \int_a^b u'(x)v(x)\, dx + \int_a^b u(x)v'(x)\, dx.
\]
Equivalent à :
\[
[u(x)v(x)]_a^b = \int_a^b u'(x)v(x)\, dx + \int_a^b u(x)v'(x)\, dx.
\]
En réarrangeant les termes, on a :
\[
\int_a^b u(x)v'(x)\, dx = [u(x)v(x)]_a^b - \int_a^b u'(x)v(x)\, dx.
\]
\end{proof}

\begin{example}
    Trouver la primitive de la fonction de logarithme naturelle \(f(x) = \ln(x)\) pour \(x > 0\).
    \noindent
    On utilise l'intégration par parties en posant :
    \[u = \ln(x) \quad \Rightarrow \quad u' = \frac{1}{x},\]
    \[v' = 1 \quad \Rightarrow \quad v = x.\]
    \noindent
    En appliquant la formule d'intégration par parties :
    \[\int \ln(x)\, dx = [\ln(x) \cdot x] - \int \frac{1}{x} \cdot x\, dx.\]
    \noindent
    Cela donne :
    \[\int \ln(x)\, dx = x \ln(x) - \int 1\, dx = x \ln(x) - x + C,\]
    où \(C\) est la constante d'intégration. 
\end{example}
\medskip

\section{\textbf{Extension aux fonctions complexes}}

Si \(f : I \to \mathbb{C}\) est une fonction complexe définie sur un intervalle \(I \subset \mathbb{R}\), on peut écrire :
\[
f(x) = u(x) + i\, v(x),
\]
où \(u,v : I \to \mathbb{R}\) sont les parties réelle et imaginaire de \(f\).  
Alors les résultats sur les intégrales de fonctions réelles s'étendent naturellement aux fonctions complexes, car :
\[
\int_a^b f(x)\, dx = \int_a^b u(x)\, dx + i \int_a^b v(x)\, dx,
\]
et toutes les propriétés précédentes (linéarité, intégration par parties, changement de variable) s'appliquent séparément aux parties réelle et imaginaire.



\end{document}